\documentclass[10pt,a4paper]{article}
\usepackage[T1]{fontenc}
\usepackage{lmodern}
\usepackage[margin=23mm]{geometry}
\usepackage{amsmath,amssymb,amsthm,mathtools}
\usepackage[hidelinks]{hyperref}
\newtheorem{theorem}{Theorem}
\newtheorem{lemma}{Lemma}
\newtheorem{proposition}{Proposition}
\theoremstyle{definition}
\newtheorem{definition}{Definition}
\theoremstyle{remark}
\newtheorem{remark}{Remark}

% !TEX root = root.tex
% Define commands for sets and operators
\newcommand{\R}{\mathbb{R}}

\newcommand{\bmat}[1]{\begin{bmatrix}#1\end{bmatrix}}
\newcommand{\ip}[2]{\left\langle #1, #2 \right\rangle}
\newcommand{\mcl}[1]{\mathcal{#1}}
\newcommand{\mbf}[1]{\mathbf{#1}}
\newcommand{\mrm}[1]{{\mathrm{#1}}}
\newcommand{\RL}[0]{\R L_2}

% Partitioned representation of a 4-PI operator:
% \PIop{P}{Q_1}{Q_2}{\{R_0,R_1,R_2\}}
\newcommand{\PIop}[4]{%
    \Pi\!\left[
    \begin{array}{c|c}
        #1 & #2\\
        \hline
        #3 & #4
    \end{array}
    \right]}

    \newcommand{\Ss}{\mathbb S}
\newcommand{\COP}{\operatorname{COP}}
\newcommand{\col}{\operatorname{col}}
\newcommand{\diag}{\operatorname{diag}}
\newcommand{\conv}{\operatorname{conv}}
\newcommand{\cone}{\operatorname{cone}}
\newcommand{\sat}{\operatorname{sat}}
\newcommand{\mc}{\mathcal M_c}
\newcommand{\Fenv}{\mathcal F_{\underline\Sigma,\overline\Sigma}}

\allowdisplaybreaks
\setlength{\parindent}{0pt}
\setlength{\parskip}{5pt}
\title{(Incremental) Quadratic Constraints for Saturations}
\author{}
\begin{document}
\maketitle
In this document we will examine \cite{noori2024completesetquadraticconstraints}
to find suitable Quadratic Constraints (QCs) and incremental QCs parameterizations
for saturations and sigmoids. The goal is eventually to connect it to \cite{lenssen2026dual} to perform 
state based controller and observer synthesis to help create a Deep Brain/Optogenetics Stimulation 
in patients suffering from Parkinsons Disease (PD). The controller consisting of a state-feedback
and Luenberger observer will then be tested on the digital twin developped in \cite{fleming2020}.

\subsection{State-feedback synthesis under input saturation verifying incremental stability}
Examine the following system
\begin{align}\label{eqn:satsys}
    T\dot{x} &= Ax + B\mathrm{sat}_{[\ell,h]}(u)\\
    u &= Kx,
\end{align}
where
\begin{equation*}
    \mathrm{sat}_{[\ell,h]}(v) := \left\{\begin{aligned}
        h \quad &\text{if} \quad v > h\\
        v \quad &\text{if} \quad \ell \leq v \leq h\\
        \ell \quad &\text{if} \quad v < \ell.
    \end{aligned}\right.
\end{equation*}
Also define the ReLU function
\begin{equation*}
    \Phi(v) := \left\{\begin{aligned}
        v \quad &\text{if} \quad v\geq 0\\
        0 \quad &\text{if} \quad v < 0.
    \end{aligned}\right.
\end{equation*}
Indeed the saturation function can be described as a combination of ReLU functions
\begin{equation*}
    \mathrm{sat}_{[\ell, h]}(v) = \ell + \Phi(v-\ell) - \Phi(v-h).
\end{equation*}
We can then describe \eqref{eqn:satsys} as
\begin{align*}
    T\dot{x} &= Ax + B\ell + B\bmat{1&-1}w \\
    w &= \Phi(z), \\
    z &= \bmat{Kx-\ell\\Kx-h}.
\end{align*}
Use the definition of incremental QCs
\begin{equation}
    \ip{\bmat{z_1-z_2\\ w_1-w_2}}{M\bmat{z_1-{z}_2\\w_1-{w}_2}}\geq 0
\end{equation}
Then introduce $x_{\delta}(t) = x_1(t) - {x}_2(t)$, ${w}_{\delta}(t) = w_1(t) - {w}_2(t)$, $z_\delta(t) = z_1(t) - {z}_2(t)$ and ${u}_{\delta}(t) = u_1(t)-{u}_2(t)$.
We gather
\begin{align}\label{eqn:lfrsatsys}
    T\dot{x}_\delta &= Ax_\delta + B\bmat{1&-1}{w}_\delta\\
    {w}_{\delta} &= \Phi(z_1)-\Phi(z_2),\\
    {z}_\delta &=\bmat{Kx_\delta\\Kx_\delta}
\end{align}
\subsection{Sigmoid with output asymptotes \(\ell<h\)}
Let \(b=(h+\ell)/2\) be the output midpoint, \(a=(h-\ell)/2\) the
half-range, \(c\) the input midpoint, and \(k>0\) the slope there.
Define \(r=k(q-c)/a\). The increasing sigmoid and its matched saturation are
\[
 F(q)=b+a\tanh r,\qquad S(q)=\mathrm{sat}_{[\ell,h]}(b+k(q-c)).
\]
Both pass through \((c,b)\) with slope \(k\). For \(c=b\) and \(k=1\),
\(S(q)=\mathrm{sat}_{[\ell,h]}(q)\) as defined above.
Since \(S(q)-b=a\,\mathrm{sat}_{[-1,1]}(r)\), we obtain
\[
 F(q)=b+\theta(q)(S(q)-b),\qquad
 \theta(q)=\frac{\tanh r}{\mathrm{sat}_{[-1,1]}(r)}
 \quad(r\ne0),\qquad \theta(c)=1 .
\]
The sharp ordinary gain bounds remain \(\tanh(1)\le\theta(q)\le1\).
For two inputs the incremental gain is instead the secant slope of \(F\).
Since \(F'(q)=k/\cosh^2(r)\), its sharp closed global interval is
\([0,k]\). If both inputs satisfy \(|q-c|\le L\) with \(L>0\), the
sharp closed interval is \([k/\cosh^2(kL/a),k]\).
\bibliographystyle{plain}
\bibliography{../../references}
\end{document}
